% The only appendix of this deck: the method chapter, written for the
% students.  This round made no new literature searches, so there is no
% backed-claim chapter -- and therefore no \chapterprecis, which belongs
% to backed-claim chapters only.  The claims that still need a search are
% listed in "Fortsatt arbete" and carry a % TODO(claim-audit) comment
% beside them in contents.nw.

\chapter{Att kontrollera ett påstående}
\label{app:verifiera}

Den här övningen gör, som allt undervisningsmaterial, en rad påståenden.
De flesta går att pröva vid datorn: att filen ligger kvar när
översättningen misslyckas ser du själv genom att köra programmet, och
att modulen \mintinline{python}{csv} sätter citattecken runt ett ord med
kommatecken i syns i körningen.  Men några av påståendena handlar om vad
forskningen, dokumentationen eller en standard säger, och dem kan ingen
körning avgöra:
\begin{enumerate}
  \item att du lär dig mer av att pröva uppgiften först och sedan jämföra
    ditt försök med lösningsförslaget än av att läsa förslaget direkt
    (inledningen);
  \item att modulen \mintinline{python}{csv} hanterar fält med
    kommatecken så som CSV-beskrivningen föreskriver, och att en slinga
    över filobjektet håller bara en rad i minnet åt gången
    (\cref{sec:rakna,sec:fortune});
  \item att nybörjare tar det som läses ur en fil för tal fastän det är
    strängar (\cref{sec:rakna});
  \item att DRY och SRP är etablerade principer med de upphov materialet
    anger (\cref{sec:filnamnet});
  \item att en JPEG-fil alltid inleds med byten \texttt{FF D8 FF}
    (\cref{sec:bild}), och att rövarspråket kommer ur Astrid Lindgrens
    \emph{Mästerdetektiven Blomkvist} (\cref{sec:rovarspraket});
  \item att det säkraste sättet att pröva ett filformat är att skriva
    och läsa tillbaka (\cref{sec:rakna}).
\end{enumerate}
Hur vet man om sådant stämmer?  Samma fråga möter dig varje gång du
läser en lärobok, en webbsida eller ett svar från en språkmodell --- och
varje gång du själv ska skriva en rapport.  Den här bilagan sammanfattar
arbetsgången; \cref{tab:paastaaenden-ovning-files} visar var svaren
redan står, och \cref{sec:ovning-files-fortsatt} vad som återstår.
Arbetsgången beskrivs utförligare i bilagan med samma namn i föreläsningen
\emph{Algoritmiskt tänkande}, och de kontroller den här övningen lutar
sig mot är gjorda i föreläsningarnas bilagor.

\begin{table}[!htb]
  \centering
  \caption{Övningens påståenden som frågor, och var svaret står.
    \enquote{Bilaga} avser bilagan i den föreläsning som namnges;
    \cref{sec:ovning-files-fortsatt} tar upp de tre frågor som ännu
    saknar svar.}
  \label{tab:paastaaenden-ovning-files}
  \begin{tabular}{@{}p{0.48\linewidth}p{0.42\linewidth}@{}}
    \toprule
    Fråga & Var svaret står \\
    \midrule
    Lär man sig mer av att lösa problemet före genomgången, och vad i
      ordningen är det som verkar? & \emph{Upprepningar}, bilaga C \\
    Gör modulen \mintinline{python}{csv} och slingan över filobjektet
      det materialet säger? & \emph{Arbeta med filer}, bilaga D och E \\
    Tar nybörjare det som läses ur en fil för tal fastän det är
      strängar? & \emph{Arbeta med filer}, bilaga B \\
    Kommer DRY från Hunt och Thomas, och är principen till nytta? &
      \emph{Algoritmiskt tänkande}, bilaga E \\
    Kommer \foreignlanguage{english}{single responsibility} från Robert
      C.\ Martin, och är principen etablerad? & \emph{Funktioner},
      bilaga B \\
    Inleds en JPEG-fil alltid med \texttt{FF D8 FF}? &
      \Cref{sec:ovning-files-fortsatt} \\
    Kommer rövarspråket ur \emph{Mästerdetektiven Blomkvist}? &
      \Cref{sec:ovning-files-fortsatt} \\
    Är att skriva och läsa tillbaka det säkraste sättet att pröva ett
      filformat? & \Cref{sec:ovning-files-fortsatt} \\
    \bottomrule
  \end{tabular}
\end{table}

\section{Gör påståendet till en fråga}

Ett påstående går att kontrollera först när det är tydligt vad som
skulle kunna visa att det är \emph{fel}.  Formulera det därför som en
fråga med ett svar.  \enquote{Sätter modulen
\mintinline{python}{csv} citattecken runt ett fält som innehåller
avgränsaren?} kan besvaras med ja eller nej; \enquote{filhantering är
viktigt} kan det inte, eftersom \enquote{viktigt} är ett omdöme och inte
ett faktum.  Ofta döljer sig flera påståenden i ett: \enquote{det är
bättre att lösa uppgiften själv först} är två frågor --- \emph{är} det
bättre, och \emph{vad} i det är det som hjälper --- och de har olika
svar.

\section{Välj rätt sorts källa}

Olika frågor besvaras av olika källor.
\begin{description}
  \item[Dokumentationen] för vad ett verktyg \emph{gör}.  Vad
    \mintinline{python}{csv.writer} gör med ett fält som innehåller ett
    kommatecken avgörs av Pythons egen dokumentation, inte av en
    lärobok som återger den.
  \item[Standarden] för vad ett format \emph{är}.  För CSV och JSON är
    det de dokument (RFC:er) som beskriver formaten; för en JPEG-fils
    inledande byte är det formatspecifikationen.
  \item[Primärkällan] för var en idé \emph{kommer ifrån}, och för ett
    specifikt resultat.  Att rövarspråket kommer ur en bestämd bok
    avgörs av boken.
  \item[Översikter] (\foreignlanguage{english}{reviews}) för om något är
    \emph{etablerat}, eller för vad forskningen \emph{sammantaget}
    säger.  Frågan om det hjälper att lösa problemet först är en sådan
    fråga, och svaret i \emph{Upprepningar}, bilaga C, vilar på en
    översikt av flera experiment.
\end{description}
Var man hittar dem: dokumentationen på \url{docs.python.org}, RFC:erna
hos \url{rfc-editor.org}, och för forskningen universitetsbibliotekets
söktjänst och databaser som Scopus, Web of Science, IEEE Xplore, DBLP
(datalogi) och OpenAlex (öppen).

\section{Sök så att du kan ha fel}

Den som söker på namn hon redan känner till hittar det hon väntar sig.
Sök därför \emph{ämnesdrivet} --- på begrepp, inte författare --- och sök
också efter \emph{invändningar}.  Hur mycket det kan betyda syns i den
här övningens inledning: den första källan säger att det hjälper att
lösa problemet före genomgången, och motsökningen gav den andra, som
visar att det inte är dröjsmålet som hjälper utan jämförelsen.  Utan
motsökningen hade materialet uppmanat dig att vänta med
lösningsförslaget, i stället för att jämföra ditt försök med det.

\section{Läs i fulltext, och skriv ner hur}

En träff är inte ett belägg.  Läs källan, citera med sidnummer, och skriv
ner hur du hittade den, varför du valde just den, vad den säger
ordagrant, och vad motsökningen gav.  I det här materialet står det i
filen \texttt{ltnotes.bib}, ett stycke per källa, och den som vill
kontrollera ett påstående kan följa spåret därifrån.  Ingen av källorna i
den här övningen söktes upp här: alla är återanvända ur
föreläsningarnas bilagor, och varje post säger vilken bilaga som besvarar
frågan.

\section{Dra slutsatsen --- med förbehåll}

Skriv påståendet så starkt som källorna bär, och inte starkare.  Därför
står det i inledningen att det hjälper \emph{måttligt} att lösa
problemet först, och att det är jämförelsen som verkar --- inte att
metoden alltid är bäst.

\section{En anmärkning om språkmodeller}

En språkmodell kan hjälpa dig att formulera frågan och att hitta
sökord, men den är ingen källa.  Den kan ange en bok, ett sidnummer och
ett citat som inte finns, och den gör det med samma säkra ton som när
den har rätt.  Kontrollera alltid i originalet --- och skriv ner att du
gjort det.

\section{Fortsatt arbete}
\label{sec:ovning-files-fortsatt}

Tre av övningens påståenden är ännu inte kontrollerade, och alla tre är
märkta med en kommentar i källan (\texttt{contents.nw}).

Det första är att en JPEG-fil alltid inleds med byten \texttt{FF D8 FF}
(\cref{sec:bild}).  Frågan är en formatfråga och ska avgöras av
formatspecifikationen --- ITU-T T.81 eller JFIF-specifikationen --- inte
av en webbsida med en tabell över \enquote{magiska tal}.  Den bör också
delas i två: inleds \emph{varje} JPEG-fil med de tre byten, och räcker de
tre för att avgöra att filen är en JPEG?  Det senare är svagare än det
förra.

Det andra är att rövarspråket kommer ur Astrid Lindgrens
\emph{Mästerdetektiven Blomkvist} (\cref{sec:rovarspraket}).  Frågan är
en fråga om upphov och ska avgöras av boken, med förbehållet att en
barnkammarkod kan ha funnits före den bok som gjorde den känd; en
sekundärkälla som säger \enquote{Astrid Lindgren populariserade} svarar
alltså på en annan fråga än \enquote{Astrid Lindgren uppfann}.

Det tredje är svagare men hör hit: att det säkraste sättet att pröva ett
filformat är att skriva och läsa tillbaka (\cref{sec:rakna}).  Det är
ärvt från föreläsningen \emph{Arbeta med filer}, där det också står utan
källa, och det bör kontrolleras där --- frågan är vad \enquote{säkraste}
betyder, och om en jämförelse med det skrivna verkligen fångar de fel
ett format kan ha.
